\documentclass[11pt,a4paper]{article}

\usepackage{algorithm}
\usepackage{algpseudocode}
\usepackage[utf8]{inputenc}
\usepackage[T1]{fontenc}
\usepackage{lmodern}
\usepackage[margin=1in]{geometry}
\usepackage{fancyhdr}
\usepackage{lastpage}
\usepackage{setspace}
\usepackage{xcolor}
\usepackage{graphicx}
\usepackage{amsmath,amssymb}
\usepackage{booktabs}
\usepackage{float}
\usepackage{pgfplots}
\pgfplotsset{compat=1.18}
\usepackage[hidelinks]{hyperref}

% ---------------------------------------------------------------
% Header / footer
% ---------------------------------------------------------------
\pagestyle{fancy}
\fancyhf{}
\fancyhead[L]{\itshape TARI29 Artificial Intelligence, 2026}
\fancyhead[R]{\itshape Evolutionary Computation Assignment}
\fancyfoot[C]{Page \thepage{} of \pageref{LastPage}}
\renewcommand{\headrulewidth}{0.4pt}
\renewcommand{\footrulewidth}{0pt}

\fancypagestyle{plain}{%
  \fancyhf{}%
  \fancyhead[L]{\itshape TARI29 Artificial Intelligence, 2026}%
  \fancyfoot[C]{Page \thepage{} of \pageref{LastPage}}%
  \renewcommand{\headrulewidth}{0.4pt}%
  \renewcommand{\footrulewidth}{0pt}%
}

% ---------------------------------------------------------------
% Title block
% ---------------------------------------------------------------
\newcommand{\groupnumber}{9}
\title{Group \groupnumber{} Assignment Report\\[0.2em]
  \large An Evolutionary Algorithm for the Travelling Salesman Problem}
\author{%
  Christos Zampounis \quad Mehmet Cihan Yalcin \quad P\'eter D\'avid Hajdu\\[0.3em]
  Pierre Kasparian \quad Prajval Parthiban}
\date{September 30, 2026}

% Red helper for the template's instructions
\newcommand{\instruction}[1]{\textcolor{red}{#1}}

\begin{document}

\maketitle
\thispagestyle{plain}

% ===============================================================
\section{Introduction}
The goal of this assignment is to design, implement and experimentally evaluate
an evolutionary algorithm for a combinatorial optimisation problem. We chose the
Travelling Salesman Problem, whose discrete, permutation-based search space and
relevance to routing and logistics make it a natural and challenging testbed for
evolutionary computation. This report presents our problem formulation, the
evolutionary approach and operators we developed, and an experimental analysis
of its solution quality and convergence behaviour.

% ===============================================================
\section{The Travelling Salesman Problem}

\subsection{Mathematical Formulation}

The Travelling Salesman Problem (TSP) is a classical combinatorial
optimisation problem that, despite its simple formulation, is difficult to
solve: given a set of $N$ cities and the distances between them, the task is
to find the shortest closed tour that visits every city exactly once. The
problem is NP-hard and no exact algorithm can solve it in polynomial time.
Following the formulation of Potvin~\cite{potvin1996}, we model the problem
on a complete weighted graph $G=(V,E)$, where $V=\{1,\dots,N\}$ is the set of
cities and $d_{ij}$ is the distance between cities $i$ and $j$. The decision
variables are $x_{ij}\in\{0,1\}$, where $x_{ij}=1$ if and only if the tour
travels directly from city $i$ to city $j$. The objective is to minimise the
total travelled distance,
\begin{equation}
  \text{Minimise} \quad \sum_{i=1}^{N}\sum_{j=1}^{N} d_{ij}\,x_{ij},
  \label{eq:tsp_objective}
\end{equation}
subject to the degree constraints
\begin{equation}
  \sum_{j=1}^{N} x_{ij}=1, \quad i=1,\dots,N,
  \qquad
  \sum_{i=1}^{N} x_{ij}=1, \quad j=1,\dots,N,
  \label{eq:tsp_degree}
\end{equation}
which require every city to have exactly one incoming and one outgoing edge,
and the subtour-elimination constraints
\begin{equation}
  \sum_{i,j\in S} x_{ij} \;\le\; |S|-1,
  \qquad \forall\, S\subset V,\; 2\le |S|\le N-1,
  \label{eq:tsp_subtour}
\end{equation}
which forbid partial cycles over any proper subset $S$ of the cities. Without
constraints~\eqref{eq:tsp_subtour}, the model reduces to the assignment problem
and its optimal solution generically consists of several disconnected subtours
rather than a single tour. For symmetric instances, such as the Euclidean TSP
(ETSP) considered in this report, the optimal solution of the assignment-problem
relaxation typically contains many two-vertex subtours; such instances are
therefore better tackled by specialised algorithms that exploit this
structure~\cite{potvin1996}.

Because an evolutionary algorithm operates on orderings rather than on binary
edge variables, it is convenient to state the equivalent tour-as-permutation
formulation. A candidate solution is a permutation
$\pi=(\pi_1,\dots,\pi_N)$ of the cities, and its cost is the length of the
cycle that visits them in that order,
\begin{equation}
  C(\pi) \;=\; \sum_{k=1}^{N-1} d_{\pi_k,\pi_{k+1}} \;+\; d_{\pi_N,\pi_1}.
  \label{eq:tsp_permutation}
\end{equation}
The evolutionary algorithm therefore seeks a permutation minimising
Equation~\eqref{eq:tsp_permutation}; this representation is the one adopted
throughout this report.

% ===============================================================
\subsection{Related Evolutionary Approaches}
\label{sec:related}

\subsubsection{Edge-Swapping Genetic Algorithm (Liu, 2014)}
Liu~\cite{liu2014} builds a genetic algorithm around a crossover operator
called \emph{Edge Swapping} (ES), which keeps the edges shared by the two
parent tours and swaps in the cycles where they disagree. The population is
initialised not randomly but with a greedy local search using the 2-opt
neighbourhood, and on large benchmarks the algorithm reaches best-known
solutions. Following Liu, we therefore seed our own population with greedy
nearest-neighbour tours refined by 2-opt.

\subsubsection{Genetic-based Particle Swarm Optimization (Liao et al., 2012)}
Liao et al.~\cite{liao2012} propose GPSO, a genetic-based particle swarm
optimisation in which each particle is a real-valued vector decoded into a tour
by ranking its components. We adopt their geometric initialisation: cities are
sorted by their angle around the geographical centre, producing a non-crossing
route. 

\subsubsection{Genetic Algorithm with Order Crossover (Hussain et al., 2017)}
Hussain et al.~\cite{hussain2017} compare crossover operators for the path
representation, where a chromosome is the visiting order and classical
one-point or uniform crossovers cannot be used as they duplicate cities. They
review partially mapped (PMX), order (OX) and cycle crossover, and propose a
modified CX2, but no operator dominated on TSPLIB. We adopt OX: a segment of
one parent is copied into the child and the remaining cities are filled in
their order of appearance in the other parent, always yielding a valid tour.

% ===============================================================
\subsection{Motivation of the Developed Approach}

The evolutionary approach we develop follows the genetic algorithm presented in
Section~9.2 of Mitchell~\cite{mitchell1997}, adapted to the permutation
representation of the TSP; its pseudo-code is given in
Section~\ref{sec:algorithm}. Within this template, the design choices left open
concern the initialization of the population and the crossover operator, which
we align with the literature reviewed in Section~\ref{sec:related}. We
initialize the population from informed tours rather than at random, combining
greedy nearest-neighbour tours refined by 2-opt~\cite{liu2014} with geometric,
non-crossing tours~\cite{liao2012}; for recombination we adopt the order
crossover (OX)~\cite{hussain2017}. The remaining operators are left
unchanged, so that the effect of the two adapted components can be isolated in
Section~\ref{sec:experimental}.

% ===============================================================
\subsection{Related Evolutionary Approaches}
\label{sec:related}

\subsubsection{Edge-Swapping Genetic Algorithm (Liu, 2014)}
Liu~\cite{liu2014} builds a genetic algorithm around a crossover operator
called \emph{Edge Swapping} (ES), which keeps the edges shared by the two
parent tours and swaps in the cycles where they disagree. The population is
initialised not randomly but with a greedy local search using the 2-opt
neighbourhood, and on large benchmarks the algorithm reaches best-known
solutions. Following Liu, we therefore seed our own population with greedy
nearest-neighbour tours refined by 2-opt.

\subsubsection{Genetic-based Particle Swarm Optimization (Liao et al., 2012)}
Liao et al.~\cite{liao2012} propose GPSO, a genetic-based particle swarm
optimisation in which each particle is a real-valued vector decoded into a tour
by ranking its components. We adopt their geometric initialisation: cities are
sorted by their angle around the geographical centre, producing a non-crossing
route. 

\subsubsection{Genetic Algorithm with Order Crossover (Hussain et al., 2017)}
Hussain et al.~\cite{hussain2017} compare crossover operators for the path
representation, where a chromosome is the visiting order and classical
one-point or uniform crossovers cannot be used as they duplicate cities. They
review partially mapped (PMX), order (OX) and cycle crossover, and propose a
modified CX2, but no operator dominated on TSPLIB. We adopt OX: a segment of
one parent is copied into the child and the remaining cities are filled in
their order of appearance in the other parent, always yielding a valid tour.

% ===============================================================
\subsection{Motivation of the Developed Approach}

The evolutionary approach we develop follows the genetic algorithm presented in
Section~9.2 of Mitchell~\cite{mitchell1997}, adapted to the permutation
representation of the TSP; its pseudo-code is given in
Section~\ref{sec:algorithm}. Within this template, the design choices left open
concern the initialization of the population and the crossover operator, which
we align with the literature reviewed in Section~\ref{sec:related}. We
initialize the population from informed tours rather than at random, combining
greedy nearest-neighbour tours refined by 2-opt~\cite{liu2014} with geometric,
non-crossing tours~\cite{liao2012}; for recombination we adopt the order
crossover (OX)~\cite{hussain2017}. The remaining operators are left
unchanged, so that the effect of the two adapted components can be isolated in
Section~\ref{sec:experimental}.

% ===============================================================
\newpage
\section{Algorithm}
\label{sec:algorithm}
Each individual is represented as a permutation of all cities, where the order
of the cities defines the tour and the final city is connected back to the first.
The initial population consists of $p$ randomly generated permutations.
The cost of an individual is given by Equation~\eqref{eq:tsp_permutation}.
Since the TSP is a minimisation problem, this cost is converted into a
maximisation-based fitness function defined as
\begin{equation}
    F(\pi) = \frac{1}{C(\pi)}.
    \label{eq:fitness}
\end{equation}
Consequently, shorter tours receive higher fitness values and are favoured
during selection.

In each generation, a fraction $(1-r)$ of the population is selected and carried over unchanged. The remaining fraction $r$ is replaced by offspring generated
from selected parent pairs. Tournament selection is used to select parents. For each parent a small number of individuals are sampled randomly from the 
population, and the individual with the highest fitness is selected. This process is repeated to obtain the second parent. Tournament selection was chosen because it
is simple to implement and provides selection pressure towards shorter tours while maintaining population diversity.

Order Crossover (OX) is used to generate offspring because it preserves valid permutations. A segment from one parent is copied to the offspring,
while the remaining cities are inserted according to their order in the second parent, skipping already used cities. Two offspring are generated for each parent pair
by swapping the parent roles.

After crossover, swap mutation is applied to a fraction $m$ of the new population. Two randomly selected positions in an individual are exchanged, which introduces variation
while preserving a valid permutation. The process is repeated for $G$ generations, and the best tour found during the run is returned.


\begin{algorithm}[H]
\caption{Genetic Algorithm for the Travelling Salesman Problem}
\begin{algorithmic}
\Require cities, population size $p$, replacement rate $r$, mutation rate $m$, number of generations $G$
\Ensure best tour found

\State $P \gets p$ random permutations of the cities
\State evaluate the fitness of each individual in $P$

\For{$g = 1$ to $G$}
    \State $P_s \gets$ the best $(1-r)p$ individuals from $P$

    \For{$i = 1$ to $(rp)/2$}
        \State select parents $h_1,h_2$ from $P$ using tournament selection
        \State $o_1 \gets \textsc{OrderCrossover}(h_1,h_2)$
        \State $o_2 \gets \textsc{OrderCrossover}(h_2,h_1)$
        \State add $o_1$ and $o_2$ to $P_s$
    \EndFor

    \State apply \textsc{SwapMutation} to a fraction $m$ of $P_s$
    \State $P \gets P_s$
    \State evaluate the fitness of each individual in $P$
    \State update the best tour found
\EndFor

\State \Return best tour

\end{algorithmic}
\end{algorithm}
% ===============================================================
\section{Experimental part}
\label{sec:experimental}

% ===============================================================
\section{Results and Analysis}

% ===============================================================
\section{Conclusions}

% ===============================================================
\begin{thebibliography}{9}

\bibitem{potvin1996}
J.-Y. Potvin, ``Genetic algorithms for the traveling salesman problem,''
\emph{Annals of Operations Research}, vol.~63, no.~3, pp.~339--370, 1996.
\url{https://doi.org/10.1007/BF02125403}.

\bibitem{liu2014}
S.~Liu, ``A Powerful Genetic Algorithm for Traveling Salesman Problem,''
arXiv:1402.4699, 2014. \url{https://arxiv.org/abs/1402.4699}.

\bibitem{liao2012}
Y.-F.~Liao, D.-H.~Yau, and C.-L.~Chen, ``Evolutionary algorithm to traveling
salesman problems,'' \emph{Computers \& Mathematics with Applications},
vol.~64, no.~5, pp.~788--797, 2012.
\url{https://doi.org/10.1016/j.camwa.2011.12.018}.

\bibitem{hussain2017}
A.~Hussain, Y.~S.~Muhammad, M.~N.~Sajid, I.~Hussain, A.~M.~Shoukry, and
S.~Gani, ``Genetic algorithm for traveling salesman problem with modified cycle
crossover operator,'' \emph{Computational Intelligence and Neuroscience},
vol.~2017, Article ID 7430125, 7 pages, 2017.
\url{https://doi.org/10.1155/2017/7430125}.

\bibitem{mitchell1997}
T.~M. Mitchell, \emph{Machine Learning}. New York, NY, USA: McGraw-Hill, 1997,
ch.~9, sec.~9.2.
\url{https://www.cs.cmu.edu/~tom/files/MachineLearningTomMitchell.pdf}.

\end{thebibliography}

\end{document}
